\documentclass[10pt,a4paper]{amsart}
\usepackage[margin=19mm]{geometry}
\usepackage{amsmath,amssymb,mathtools}
\usepackage[scr=rsfs,cal=euler]{mathalfa}
\usepackage{xfrac}
\usepackage[unicode,hidelinks,bookmarks=false]{hyperref}
\newcommand{\ie}{i.e.\ }
\newcommand{\cf}{cf.\ }
\newcommand{\resp}{respectively\ }
\newcommand{\Subsection}{\subsection}
\providecommand{\nobreakdash}{\nobreak\hbox{-}}
\setlength{\emergencystretch}{2em}
\multlinegap0pt
\usepackage{bm}
\usepackage{bbm}
\usepackage{dsfont}
\usepackage{xspace}
\usepackage{cancel}
\usepackage{mathtools}
\usepackage{amsthm}
\makeatletter
\@ifundefined{lemma}{\newtheorem{lemma}{Lemma}}{}
\makeatother
\title[Policy Optimization and Statistical Inference for Online Con]{Policy Optimization and Statistical Inference for Online Contextual Matrix Games}
\author{Liner Xiang, Yixin Wang, Hengrui Cai}
\date{Source: arXiv:2608.17173v1 (CC BY 4.0)}
\begin{document}
\maketitle

\noindent\emph{Source and licence.} Policy Optimization and Statistical Inference for Online Contextual Matrix Games. Version arXiv:2608.17173v1 (CC BY 4.0), distributed under the Creative Commons Attribution 4.0 International licence, as recorded in the arXiv licence metadata for this exact version and shown on the abstract page https://arxiv.org/abs/2608.17173v1. Full rights notice: licenses/arxiv-2608.17173-CC-BY-4.0.txt.\par\smallskip

\noindent\textit{Editorial scope.} Counted unit: the Lemma whose source label is \texttt{lem:context\_ipw\_score\_concentration} in the article (source0.tex lines 5745--5771, source label \texttt{lem:context\_ipw\_score\_concentration}), delivered together with the article's own complete original proof of it (source0.tex lines 5773--6100, one \texttt{proof} environment of 328 lines). No mathematics is written, repaired or re-derived: statement and proof are the article's own text line for line. Required context retained uncounted: none. Every definition, display and label of those lines is retained unchanged. Omitted from this package and not redistributed: 1--5744, 5772--5772, 6101--6131. Nothing inside a retained excerpt is omitted or altered: every retained body is delivered exactly as the article's lines are written, so no omission receipt of any kind accompanies this package. Citations: the retained excerpts carry no citation command at all, so no bibliography entry of the article is transcribed and no reference is redistributed. Reference adaptation: none was needed and none was made; every reference inside the delivered unit resolves inside it, so no unresolved reference or citation is produced. Printed slips kept literal: no printed word, symbol, hypothesis or number of the counted statement or of its proof was altered. Layout and class: the article's own class is \texttt{article} with packages including amsmath, graphicx, psfrag, epsf, color, enumerate, natbib, subcaption, amsthm, mathtools, amsfonts, bm, hyperref, setspace; the wrapper is the lane's pinned \texttt{amsart} editorial wrapper at 10pt on A4, which restates the theorem-like environments the retained text needs and adds the article's own macro definitions that the retained text needs (none), with extra wrapper packages (bm, bbm, dsfont, xspace, cancel, mathtools). The delivered numbering is the unit's own. Assets: no retained span contains a figure environment, a graphics-inclusion command, a PDF-page inclusion, a table or a TikZ picture; no image file is copied into this package, the package declares no asset dependency and no image file is redistributed with it. Licence: version arXiv:2608.17173v1 of this article is distributed under the Creative Commons Attribution 4.0 International licence, as recorded in the arXiv licence metadata for this exact version and shown on the abstract page https://arxiv.org/abs/2608.17173v1. A full rights notice is delivered at licenses/arxiv-2608.17173-CC-BY-4.0.txt. Characters: every retained excerpt is pure ASCII, so the delivered engine, XeLaTeX with its default Latin Modern text font, reports no missing character.
\begin{lemma}[Concentration of the contextual IPW terms]
\label{lem:context_ipw_score_concentration}
Fix an action pair $(i,j)$. Let $a_s^{ij}=\mathbb I\{i_s=i,j_s=j\}$ and $\pi_s(i,j)=\mathbb P\left(i_s=i,j_s=j\mid\mathcal H_{s-1}^{\mathrm{con}},\bm z_s\right)$.
Suppose that $\epsilon_s$ is non-increasing and $\pi_s(i,j)\ge {\epsilon_s^2}/{mk}$ for all $s\le t$.
Suppose that $\|\bm z\|_2 \le L_z$, $|\varphi^{ij}(\bm z) - \bm z^\top\bm\beta_\dagger^{ij}| \le L_\varphi$ for any $\bm{z} \sim P_z$. 
Then, for any $u>0$ and some $\ell \in [d]$,
\begin{equation}
\label{equ:ipw_approximation_score_tail}
\mathbb P
\left(\left|
\frac{1}{t}\sum_{s=1}^t\frac{a_s^{ij}}{\pi_s(i,j)}z_{s,\ell} \left(\varphi^{ij}(\bm z_s) - \bm z_s^\top\bm\beta_\dagger^{ij}\right)\right|> u \right)
\le
2\exp\left\{-\frac{t\epsilon_t^2u^2}{2mk\left(L_z^2L_\varphi^2+L_zL_\varphi u/3\right)}\right\}.    
\end{equation}
Suppose additionally that, whenever $(i_s,j_s)=(i,j)$, the reward noise $e_s$ is conditionally mean-zero and $\sigma_{ij}$-sub-Gaussian. Then, for any $u>0$ and some $\ell \in [d]$,
\begin{equation}
\label{equ:ipw_noise_score_tail}
\mathbb P
\left(\left|
\frac{1}{t}\sum_{s=1}^t\frac{a_s^{ij}}{\pi_s(i,j)}z_{s,\ell}e_s
\right|>u\right)
\le
2\exp\left[-\frac{t\epsilon_t^2}{4mk}\min\left\{\frac{u^2}{L_z^2\sigma_{ij}^2},\frac{u}{L_z\sigma_{ij}}
\right\}
\right].
\end{equation}
\end{lemma}

\begin{proof}
We first prove~\eqref{equ:ipw_approximation_score_tail}. 
By iterated conditional expectation,
\begin{align*}
& \mathbb E
\left[
\frac{ a_s^{ij} }{ \pi_s(i,j) } z_{s,\ell} \left(\varphi^{ij}(\bm z_s) - \bm z_s^\top\bm\beta_\dagger^{ij} \right)
\mid
\mathcal H_{s-1}^{\mathrm{con}}
\right] \\
=&
\mathbb E
\left[\left.
\mathbb E\left(
\left.\frac{a_s^{ij}}{\pi_s(i,j)}z_{s,\ell}
\left(\varphi^{ij}(\bm z_s) - \bm z_s^\top\bm\beta_\dagger^{ij} \right)
\right|
\mathcal H_{s-1}^{\mathrm{con}},\bm z_s
\right)\right|\mathcal H_{s-1}^{\mathrm{con}}\right] \\
=&
\mathbb E
\left[
z_{s,\ell} \left.
\left(\varphi^{ij}(\bm z_s) - \bm z_s^\top\bm\beta_\dagger^{ij} \right) \right|\mathcal H_{s-1}^{\mathrm{con}} 
\right] \\
=& 0,
\end{align*}
where the last equality follows from the first-order optimality condition for the least-false parameter. Hence, $\{\frac{ a_s^{ij} }{ \pi_s(i,j) } z_{s,\ell} \left(\varphi^{ij}(\bm z_s) - \bm z_s^\top\bm\beta_\dagger^{ij} \right)\}_{s=1}^t$ is a martingale-difference sequence.
Moreover, since $\pi_s(i,j) \ge {\epsilon_t^2}/{mk}$, we have
\[
\left|
\frac{ a_s^{ij} }{ \pi_s(i,j) } z_{s,\ell} \left(\varphi^{ij}(\bm z_s) - \bm z_s^\top\bm\beta_\dagger^{ij} \right)
\right|
\le
\frac{mk L_zL_\varphi}{\epsilon_t^2}.
\]
Its conditional variance satisfies
\small{
\begin{equation*}
\mathbb E
\left[
\left.
\left(
\frac{ a_s^{ij} }{ \pi_s(i,j) } z_{s,\ell} \left(\varphi^{ij}(\bm z_s) - \bm z_s^\top\bm\beta_\dagger^{ij} \right)
\right)^2
\right|
\mathcal H_{s-1}^{\mathrm{con}}
\right]
\le
\mathbb E
\left[
\left.
\frac{
z_{s,\ell}^2
\left\{
\left(\varphi^{ij}(\bm z_s) - \bm z_s^\top\bm\beta_\dagger^{ij} \right)
\right\}^2}{\pi_s(i,j)}
\right|
\mathcal H_{s-1}^{\mathrm{con}}
\right]
\le 
\frac{mk L_z^2L_\varphi^2}{\epsilon_t^2}.
\end{equation*}
}
Therefore,
\[
\sum_{s=1}^t
\mathbb E
\left[
\left.
\left(
\frac{ a_s^{ij} }{ \pi_s(i,j) } z_{s,\ell} \left(\varphi^{ij}(\bm z_s) - \bm z_s^\top\bm\beta_\dagger^{ij} \right)
\right)^2
\right|
\mathcal H_{s-1}^{\mathrm{con}}
\right]
\le
\frac{tmkL_z^2L_\varphi^2}{\epsilon_t^2}.
\]
Applying Freedman's inequality yields
\begin{align*}
\mathbb P
\left(
\left|
\sum_{s=1}^t
\frac{ a_s^{ij} }{ \pi_s(i,j) } z_{s,\ell} \left(\varphi^{ij}(\bm z_s) - \bm z_s^\top\bm\beta_\dagger^{ij} \right)
\right|>tu
\right)
\le & \,
2\exp
\left\{
-
\frac{t^2u^2}{2\left[tmkL_z^2L_\varphi^2/\epsilon_t^2
+ (mk L_zL_\varphi/\epsilon_t^2)tu/3
\right]
}
\right\}
\\
= & \,
2\exp\left\{
-\frac{t\epsilon_t^2u^2/mk}{2\left(L_z^2L_\varphi^2+L_zL_\varphi u/3\right)}
\right\}.
\end{align*}
This proves~\eqref{equ:ipw_approximation_score_tail}.

We next prove~\eqref{equ:ipw_noise_score_tail}. Conditional on
$\mathcal H_{s-1}^{\mathrm{con}}$ and $\bm z_s$, we have
\begin{align*}
&
\mathbb E
\left[
\left.
\exp
\left(
q\frac{ a_s^{ij} }{ \pi_s(i,j) } z_{s,\ell} e_s
\right)
\right|
\mathcal H_{s-1}^{\mathrm{con}},
\bm z_s
\right] \\
=&
1-\pi_s(i,j)+ \pi_s(i,j)
\mathbb E
\left[
\left.
\exp
\left\{
\frac{
qz_{s,\ell}e_s}{\pi_s(i,j)}
\right\}
\right|
\mathcal H_{s-1}^{\mathrm{con}},
\bm z_s,i_s=i,j_s=j
\right] \\
\le &
1-\pi_s(i,j)+\pi_s(i,j)
\exp
\left\{
\frac{q^2z_{s,\ell}^2\sigma_{ij}^2}{2\pi_s^2(i,j)}
\right\},
\end{align*}
where the last inequality comes from the conditional sub-Gaussianity of $e_s$. 
For $|q| \le {\epsilon_t^2}/{2mkL_z\sigma_{ij}}$,
we have
\[
\frac{q^2z_{s,\ell}^2\sigma_{ij}^2}{2\pi_s^2(i,j)}
\le
\frac{\left(\frac{\epsilon_t^2}{2mkL_z\sigma_{ij}}\right)^2L_z^2\sigma_{ij}^2}{2\left(\frac{\epsilon_t^2}{mk} \right)^2}
\le \frac18.
\]
Using $\exp(v)-1 \le 2v$ when $0\le v\le\frac18$, we obtain
\begin{equation*}
\mathbb E
\left[
\left.
\exp
\left(
q\frac{ a_s^{ij} }{ \pi_s(i,j) } z_{s,\ell} e_s
\right)
\right|
\mathcal H_{s-1}^{\mathrm{con}},
\bm z_s
\right]
\le
1+\frac{q^2z_{s,\ell}^2\sigma_{ij}^2}{\pi_s(i,j)}
\le
\exp
\left\{
\frac{mk q^2L_z^2\sigma_{ij}^2}{\epsilon_s^2}
\right\}.
\end{equation*}
Iterating this conditional moment-generating-function bound gives
\[
\mathbb E
\left[
\exp
\left(
q\sum_{s=1}^t
\frac{ a_s^{ij} }{ \pi_s(i,j) } z_{s,\ell} e_s
\right)
\right]
\le
\exp
\left\{
\frac{tmkq^2L_z^2\sigma_{ij}^2}{\epsilon_t^2}
\right\}.
\]
We now apply the Chernoff method. Fix any $0<q\le {\epsilon_t^2}/{2mkL_z\sigma_{ij}}$.
Since the exponential function is strictly increasing, we have
\begin{equation*}
\mathbb P
\left(
\frac{1}{t}
\sum_{s=1}^t
\frac{ a_s^{ij} }{ \pi_s(i,j) } z_{s,\ell} e_s
>u\right)=
\mathbb P
\left(
\exp
\left\{q
\sum_{s=1}^t
\frac{ a_s^{ij} }{ \pi_s(i,j) } z_{s,\ell} e_s
\right\}
>
\exp\left\{qtu\right\}\right).
\end{equation*}
Applying Markov's inequality to the nonnegative random variable $\exp
\left\{
q\sum_{s=1}^t
\frac{ a_s^{ij} }{ \pi_s(i,j) } z_{s,\ell} e_s
\right\}$,
we obtain
\begin{align*}
\mathbb P
\left(
\exp
\left\{
q\sum_{s=1}^t
\frac{ a_s^{ij} }{ \pi_s(i,j) } z_{s,\ell} e_s
\right\}
>
\exp
\left\{
qtu
\right\}
\right)
&\le
\exp
\left\{-qtu\right\}
\mathbb E
\left[
\exp
\left\{
q\sum_{s=1}^t
\frac{ a_s^{ij} }{ \pi_s(i,j) } z_{s,\ell} e_s
\right\}
\right].
\end{align*}
Using the moment-generating-function bound derived above, we further have
\begin{align*}
\exp
\left\{
-qtu
\right\}
\mathbb E
\left[
\exp
\left\{
q\sum_{s=1}^t
\frac{ a_s^{ij} }{ \pi_s(i,j) } z_{s,\ell} e_s
\right\}
\right]
\le
\exp
\left\{-qtu+\frac{tmkq^2L_z^2\sigma_{ij}^2}{\epsilon_t^2}
\right\}.
\end{align*}
Therefore, for every $0<q\le\frac{\underline\pi_t}{2L_z\sigma_{ij}}$, we obtain
\begin{equation}
\label{equ:ipw_noise_chernoff}
\mathbb P
\left(
\frac{1}{t}
\sum_{s=1}^t
\frac{ a_s^{ij} }{ \pi_s(i,j) } z_{s,\ell} e_s
>u
\right)
\le
\exp
\left\{
-qtu+\frac{tmkq^2L_z^2\sigma_{ij}^2}{\epsilon_t^2}
\right\}.
\end{equation}
The unconstrained minimizer of the exponent is $q^\ast =\frac{u\epsilon_t^2}{2mkL_z^2\sigma_{ij}^2}$.
If $u\le L_z\sigma_{ij}$, then $q^\ast$ lies in the admissible range, and taking $q=q^\ast$ yields
\[
\mathbb P
\left(
\frac1t
\sum_{s=1}^t
\frac{ a_s^{ij} }{ \pi_s(i,j) } z_{s,\ell} e_s
>u\right)
\le
\exp
\left\{
-\frac{
t\epsilon_t^2u^2}{4mkL_z^2\sigma_{ij}^2}
\right\}.
\]
If $u>L_z\sigma_{ij}$, we take $q=\frac{\epsilon_t^2}{2mkL_z\sigma_{ij}}$, which gives
\[
\mathbb P
\left(
\frac1t
\sum_{s=1}^t
\frac{ a_s^{ij} }{ \pi_s(i,j) } z_{s,\ell} e_s
>u\right)
\le
\exp
\left\{
-\frac{t\epsilon_t^2u}{4mkL_z\sigma_{ij}}
\right\}.
\]
Combining the two cases and applying the same argument to
$-\frac{ a_s^{ij} }{ \pi_s(i,j) } z_{s,\ell} e_s$ yields
\[
\mathbb P
\left(
\left|
\frac1t
\sum_{s=1}^t
\frac{ a_s^{ij} }{ \pi_s(i,j) } z_{s,\ell} e_s
\right|
>u
\right)
\le
2\exp
\left[
-
\frac{t\epsilon_t^2}{4mk}
\min
\left\{
\frac{u^2}{L_z^2\sigma_{ij}^2},
\frac{u}{L_z\sigma_{ij}}
\right\}
\right].
\]
\end{proof}

\end{document}
